% related_work.tex — Task 9 deliverable (drafted 2026-08-17); condensed and
% re-argued 2026-09-05. The 2026-08-17 draft enumerated systems roughly one per
% sentence and ran ~940 words. Each subsection now carries a single claim and
% groups its citations under it, at ~735 words (-22%). Citation COVERAGE is
% unchanged — every key cited by the earlier draft is still cited here, and
% nothing else in the paper cites a key that this section dropped.
% Citations use natbib (\citep/\citet) against ../references.bib.
% Guardrails observed: no unmeasured Medium numbers; no confidence-gate claim
% (verification is blanket, confidence gates nothing); no invented fine-tuning
% cost comparison — the in-context argument is made on no-labeled-corpus grounds.
% One factual softening vs. the 2026-08-17 draft: aggarwal2025fiscal is credited
% with validating against the documents' "own internal consistency", not
% specifically their "arithmetic structure", which REVIEW_references.md
% (2026-09-04) could not confirm from the abstract.

\section{Related Work}
\label{sec:related}

\subsection{Hierarchy recovery from documents}

Recovering logical structure from formatted documents has a long lineage ---
scholarly-PDF pipelines~\citep{lopez2009grobid}, page-object
datasets~\citep{zhong2019publaynet}, layout-aware
encoders~\citep{xu2020layoutlm,huang2022layoutlmv3}, OCR-free
vision--encoder--decoder models~\citep{kim2022donut,blecher2023nougat} --- but
its targets are flat: classification, key--value extraction and markup
transcription rather than a tree. Work that does reconstruct the tree parses it
from renderings~\citep{rausch2021docparser}, decomposes it into page-object
detection, reading order and tree
construction~\citep{ma2023hrdoc,wang2024detect}, recasts those subtasks as
relation prediction~\citep{wang2025unihdsa}, extends them to multi-domain
corpora~\citep{xing2024dochienet}, or reads table-of-contents metadata from
deeply nested legal books~\citep{wehnert2025hips}.

These systems differ in architecture but share a premise, and it is the
premise rather than any individual design that my setting violates: hierarchy
can be \emph{learned}, because supervision exists --- line labels, bounding
boxes, TOC metadata, pretraining on labeled pages --- over a genre whose
typography is internally consistent. Early learning standards offer neither.
With every jurisdiction designing its own document~\citep{ncecqa_elgs}, there
is no shared typography to learn and no annotated corpus, at any scale, to
learn it from.
What the documents do share is not appearance but nesting: whatever a level is
called, it sits at a fixed depth relative to its neighbors. My depth-map pass
is functionally the table-of-contents subtask of \citet{wang2024detect}, but
inferred zero-shot from content rather than predicted by a trained component,
then treated as the authority against which every element is classified. The
closest zero-shot neighbor, DocsRay~\citep{jeong2025docsray}, likewise prompts
for a pseudo table of contents with no training, but targets a segmentation
guiding retrieval for question answering rather than a typed hierarchy carrying
codes and ancestry, and is not graded against structural annotations.

\subsection{In-context extraction without a labeled corpus}

Prompting alone carries many tasks~\citep{brown2020language}, including
multi-step reasoning~\citep{wei2022chain}, and generative extraction is now
competitive across classic IE
tasks~\citep{wadhwa2023revisiting,xu2024large}, notably where labeled corpora
are scarce~\citep{agrawal2022large}. GoLLIE~\citep{sainz2024gollie} is the
closest methodological analogue: fine-tuning a model to follow annotation
\emph{guidelines} makes it generalize to unseen schemas. My prompts
operationalize that insight without fine-tuning, stating document-agnostic
structural principles (e.g., a structural label and its identifier constitute
the code; the text after the colon is the title) that a frozen model applies to
unseen states. This is the available move rather than merely the cheap one: no
annotated corpus of parsed early learning standards exists, so the supervised
alternative's dominant cost would be annotation, not compute
\citep[cf.][where labels do exist]{liu2022fewshot}.

Three system choices follow: schema-validated JSON with deterministic repair
rather than constrained decoding~\citep{willard2023efficient}, which keeps the
method API-agnostic; domain-scoped chunking
(\S\ref{sec:system-architecture}) rather than whole-document prompting, given
positional degradation in long contexts~\citep{liu2024lost}; and reported
run-to-run stability (\S\ref{sec:experiments-stability}) rather than a single
run, given sampling variance at fixed prompts~\citep{wang2023selfconsistency}.
Hallucination, the residual risk of generative extraction~\citep{ji2023survey},
I quantify by auditing every unmatched detection by hand. The nearest applied
system~\citep{aggarwal2025fiscal} converts government fiscal PDFs into a
research database with a multi-stage LLM pipeline validated against the
documents' own internal consistency; standards documents afford the analogous
check on their \emph{hierarchical} consistency --- codes as namespace paths,
parent--child prefix relations.

The absence of labels also decides where human effort goes: to verification
rather than annotation. Every extracted record is human-reviewed at all four
hierarchy levels --- the model drafts, humans supply the correctness guarantees
it cannot --- following the machine-teaching paradigm of experts controlling
the process rather than labeling
data~\citep{mosqueirarey2023human,monarch2021human}. It differs from active
learning~\citep{settles2009active} and from work allocating annotation between
models and humans~\citep{gilardi2023chatgpt,li2023coannotating,tan2024large} in
that verification is blanket, not selective: nothing decides \emph{which}
records a human sees, and the detector's confidence scores are recorded but
gate nothing.

\subsection{Machine-readable education standards}

Two decades of work on standards as data share one assumption. The Achievement
Standards Network modeled US K--12 standards in
RDF~\citep{sutton2008achievement} and the CASE specification defines an
exchange format for competency frameworks~\citep{1edtech2017case}, but ASN was
populated by manual modeling and CASE is a format, not an extractor.
Downstream, alignment between standards, assessments and instruction is
measured by trained human
panels~\citep{porter2002measuring,polikoff2011howwell}, while NLP from the ASN
era onward assigns learning resources~\citep{devaul2011computer} and, more
recently, assessment items~\citep{xu2025automated} \emph{to} standards already
normalized. For early childhood the assumption fails outright: no national
framework spans the age band --- the federal ELOF covers Head Start programs
only~\citep{ohs2015elof} --- and each state publishes its own guidelines as an
independently designed PDF~\citep{ncecqa_elgs,ecs2024governance}. The extraction
step this paper automates is not a competitor to that literature but its
missing input.
